\chapter{Alpha-Fetoprotein (AFP)}

\section*{What Is This Test?}
Alpha-fetoprotein (AFP) is a glycoprotein produced by the fetal yolk sac and liver. After birth, AFP levels fall rapidly, reaching adult levels by 1 year. In adults, AFP is measured as a serum tumor marker for hepatocellular carcinoma (HCC) and germ cell tumors, and as a maternal serum marker for pregnancy screening. This chapter focuses on serum AFP as a tumor marker; maternal serum AFP for pregnancy screening is also discussed briefly. In oncology, AFP is most useful in monitoring known HCC treated with resection or ablation, where rising AFP precedes imaging evidence of recurrence in many patients. AFP should be interpreted with imaging, not alone, because 20\% of HCCs are AFP-negative.

\section*{Alternative Names}
AFP; Alpha-Fetoprotein; Alpha-1-Fetoprotein; $\alpha$FP

\section*{Principle}
AFP is measured by two-site immunoassays using two monoclonal antibodies that bind different epitopes of the AFP molecule. The most common platforms are chemiluminescent immunoassay (CLIA) on automated analyzers such as Roche cobas e601, Siemens ADVIA Centaur, and Abbott Architect. Liquid chromatography-tandem mass spectrometry is increasingly used for standardization. AFP-L3 fraction (lectin-reactive AFP) is reported separately by some laboratories to improve HCC specificity.

\section*{History}
AFP was first identified in 1956 by Bergstrand and Czar in human fetal serum \cite{bergstrand1956demonstration}. Its tumor marker role emerged in the 1960s when elevated levels were documented in patients with hepatocellular carcinoma and yolk sac tumors \cite{abelev1963production}. Wide clinical use followed development of radioimmunoassays in the 1970s. Modern automated immunoassays deliver results within 30 minutes.

\section*{Why This Test Is Ordered}
\begin{itemize}
  \item Surveillance for hepatocellular carcinoma in patients with cirrhosis, chronic hepatitis B or C, or other high-risk conditions (ultrasound $\pm$ AFP every 6 months) \cite{marrero2018diagnosis}
  \item Monitoring treatment response and detecting recurrence after resection, ablation, or transplant
  \item Diagnosis and staging of germ cell tumors (non-seminomatous germ cell tumors)
  \item Evaluation of pregnant women at 15--22 weeks for fetal neural tube defects and chromosomal abnormalities
  \item Investigation of unexplained elevated liver mass
\end{itemize}

\section*{Primary vs Secondary Test}
Secondary test; AFP is used to monitor suspected or confirmed disease rather than screening asymptomatic low-risk individuals. Surveillance in high-risk patients with cirrhosis combines ultrasound and AFP.

\section*{How This Test Is Performed}
A venous blood sample is drawn into a serum separator tube (gold-top or tiger-top). 1 mL of serum suffices. For pregnancy screening, maternal serum is drawn between 15 and 22 weeks gestation (optimal 16--18 weeks) and interpreted with ultrasound dating, gestational age, maternal weight, race, and diabetes status.

\section*{What Preparation Is Needed}
None for oncology use. For pregnancy screening, accurate gestational dating by first-trimester ultrasound is essential.

\section*{Contraindications and Precautions}
False-positive AFP elevations occur in pregnancy with incorrect dating, multiple gestation, fetal demise, miscarriage, open fetal defects; in non-pregnant adults due to chronic hepatitis, cirrhosis, hepatic regeneration. AFP-L3 improves specificity when chronic liver disease is present. Anti-recombinant mouse antibodies and heterophilic antibodies can cause spurious results.

\section*{Risks}
Minimal; venipuncture only.

\section*{What Happens After the Test}
Results are available within 1--2 days. An AFP $>$ 20 ng/mL in a cirrhotic patient warrants multiphasic CT or MRI for HCC. Serial AFP trends are more reliable than a single value. In germ cell tumors, serial AFP trends assess response to chemotherapy.

\section*{Understanding Results}

\subsection*{Below Normal}
Low or undetectable AFP is normal in healthy adults and after successful treatment. In HCC, 15--20\% of tumors do not produce AFP; thus, a normal AFP does not exclude HCC in a high-risk patient. Imaging is required.

\subsection*{Above Normal}
\begin{itemize}
  \item Hepatocellular carcinoma: serum AFP $>$ 400 ng/mL with a liver mass is essentially diagnostic; levels $>$ 200 ng/mL warrant urgent imaging
  \item Chronic liver disease: moderate elevations (10--200 ng/mL) in hepatitis, cirrhosis, hepatic regeneration
  \item Germ cell tumors: non-seminomatous germ cell tumors (yolk sac tumor, embryonal carcinoma), especially with testicular or mediastinal primary
  \item Pregnancy: elevated maternal serum AFP prompts amniocentesis, high-resolution ultrasound; etiologies include neural tube defects, omphalocele, gastroschisis, incorrect dating, multiple gestation
  \item Ataxia-telangiectasia
  \item Benign conditions: hepatic regeneration, pregnancy, rare benign tumors
\end{itemize}

\section*{Normal Results}
\begin{itemize}
  \item \textbf{Adult non-pregnant:} $<$ 7 ng/mL
  \item \textbf{Pregnancy (15--22 weeks, MOM):} 0.25--2.50 MOM (multiples of the median)
  \item \textbf{HCC interpretation:} AFP elevation should prompt risk-based clinical and imaging assessment; no AFP concentration alone diagnoses HCC. Use current multiphasic imaging/LI-RADS or pathology guidance.
  \item \textbf{Germ cell tumors post-orchiectomy:} AFP should fall with half-life of 5--7 days; persistent elevation suggests residual disease
\end{itemize}

\section*{Follow-up Tests}
\begin{itemize}
  \item \textbf{Multiphasic liver CT or MRI:} For HCC characterization (arterial enhancement and portal venous washout)
  \item \textbf{Abdominal ultrasound:} Surveillance in cirrhosis
  \item \textbf{AFP-L3\%:} Improve specificity when chronic liver disease is present; L3 $>$ 10\% more specific for HCC \cite{li2001afpl3}
  \item \textbf{DCP (PIVKA-II):} Complementary HCC marker
  \item \textbf{hCG:} Together with AFP in germ cell tumor workup; seminomas produce hCG only, never AFP
  \item \textbf{Testicular ultrasound:} If germ cell tumor suspected
  \item \textbf{Amniocentesis and acetylcholinesterase:} If maternal serum AFP elevated
\end{itemize}

\section*{References}
\bibliographystyle{unsrtnat}
\bibliography{references}

\clearpage
